\documentclass[11pt]{article}
\usepackage[T1]{fontenc}
\usepackage{lmodern}
\usepackage{amsmath,amssymb}
\usepackage{booktabs}
\usepackage{array}
\usepackage[a4paper,margin=2.5cm]{geometry}
\usepackage[hidelinks]{hyperref}
\newcommand{\doi}[1]{\href{https://doi.org/#1}{\texttt{#1}}}
\setlength{\parskip}{4pt}
\title{\vspace{-1.6cm}\Large Lamb Wave Attenuation in a Layered Magneto-Electro-Elastic Plate\\[2pt]
\large with a Relaxation-Time Heat Equation and a Nonlocal Strain Gradient Elastic Stress\\[8pt]
\large Research plan, first half}
\author{Dr.\ Vipin \\ Gurugram University}
\date{30 October 2026}
\begin{document}
\maketitle
\vspace{-6mm}

\section{Title, abstract, keywords}

\paragraph{Working title} Attenuation of Lamb waves in a BaTiO$_3$/CoFe$_2$O$_4$ plate: pyroelectric and
pyromagnetic coupling, a relaxation-time heat equation and a size effect.

\paragraph{Abstract} We study plane-strain Lamb waves in a laminated plate of barium titanate and cobalt
ferrite, with free faces at $z=\pm h$ and the layering along the thickness. The model carries five fields:
the two displacements, the electric potential, the magnetic potential and the temperature rise. The
constitutive law is the linear magneto-electro-elastic law with the thermal terms retained in all three
relations, that is, thermal stress in the mechanical relation, pyroelectric polarisation in the electric
displacement and pyromagnetic induction in the magnetic induction. The heat equation is of the modified
gradient type in its base form with one phase lag, and it carries the electrocaloric and magnetocaloric
couplings as well as the thermoelastic one. The mechanical stress carries the nonlocal strain gradient
operator on its elastic part only. The magnetoelectric coefficient is not treated as a material constant but
as an effective quantity computed by two homogenisation routes, carried through the paper as a band. The
paper asks whether the thermal couplings produce an attenuation that is larger than the spread between those
two routes, and if so over what thickness and frequency range. The plan gives the model, the sequence of
comparisons, the figures, and an honest statement of how narrow the claim is.

\paragraph{Keywords} magneto-electro-elastic plate; Lamb waves; pyroelectric; pyromagnetic; magnetoelectric
composite; attenuation.

\section{Introduction}

A magneto-electro-elastic material converts among mechanical, electric and magnetic fields. No single phase
does all three well, so the practical materials are composites of a piezoelectric such as barium titanate and
a piezomagnetic such as cobalt ferrite, and the magnetoelectric effect they show is a product property of the
two phases rather than a property of either. That is the whole reason the field exists, and it is also the
reason a paper on such a material must be careful about what its coefficients mean. A magnetoelectric
coefficient measured on a laminate is a property of the laminate and of its interface, not of the material.

Wave propagation in magneto-electro-elastic multilayered plates has been solved in closed form by Chen, Pan
and Chen \cite{Chen2007}, and guided waves in a graded magneto-electro-elastic plate have been computed with
spectral elements by Xiao, Han, Liu and Li \cite{Xiao2016}. Both papers are isothermal in the sense that the
temperature is at most a constant that shifts the coefficients. The thermal couplings themselves are
classical and well measured: a change of temperature polarises the piezoelectric phase and magnetises the
piezomagnetic one, and the coefficients are called pyroelectric and pyromagnetic. Kondaiah, Shankar and
Ganesan kept both in a magneto-electro-elastic plate and showed that they change the static and the
free-vibration response \cite{Kondaiah2013}. Tocci Monaco, Fantuzzi, Fabbrocino and Luciano kept both in a
nonlocal strain gradient magneto-electro-elastic nanoplate with hygro-thermal loads \cite{Tocci2021}. So the
statement that plate papers drop these terms is not true, and we should not write it.

What is missing is the wave attenuation problem. Two recent papers mark the edge of what exists. Kakar and
Kakar solved edge waves in an electro-magneto-thermoelastic homogeneous plate \cite{Kakar2015}, which is a
plate with heat but a homogeneous material and edge rather than Lamb modes. Matlam and Kedar applied the
modified gradient theory, with memory and nonlocality, to an electro-magneto-thermoelastic medium
\cite{Matlam2026}, which is the right heat operator but not a plate and not a layered composite. To our
knowledge, after searching Crossref and the open web in October 2026 for guided waves in magneto-electro-elastic
plates, nobody has reported the dispersion and attenuation of Lamb modes in a layered
barium titanate and cobalt ferrite plate whose heat conduction carries a relaxation time, with pyroelectric
and pyromagnetic coupling both present and with a mechanical size effect.

That is a narrow claim and the plan treats it as narrow. Two consequences follow. First, the paper cannot be
sold on the couplings, because the couplings are classical; it has to be sold on the attenuation curves and
on the thickness range where the thermal effect is resolvable. Second, the uncertainty in the effective
magnetoelectric coefficient is not a nuisance to be hidden but the main competitor of the effect being
studied, so it enters every figure as a band.

\section{Problem statement and assumptions}

The plate has thickness $2h$ and free faces at $z=\pm h$. The layering is along $z$, with the two phases
arranged in a periodic or a graded stack, and the effective constants of the stack are computed once and
used throughout. Motion is plane strain in the $x$--$z$ plane. Fields vary as $\exp[i(kx-\omega t)]$ with
$k=k_r+ik_i$, a positive $k_i$ meaning decay along $+x$. Frequency is reported through $f\,(2h)$. The five
unknowns are $u_x$, $u_z$, the electric potential $\varphi$, the magnetic potential $\psi$ with
$H_i=-\psi_{,i}$, and the temperature rise $\theta$.

\begin{enumerate}
\item Linear magneto-electro-elasticity with constant effective constants at the reference temperature.
\item The effective constants come from homogenisation of the two phases, and the spread between two routes
is carried as a band rather than resolved by preference.
\item The magnetoelectric coefficient is a product property; it is not fitted to any measurement, and when a
measured value is quoted for comparison it is quoted with its own uncertainty.
\item Small temperature rise, $|\theta|/T_0\ll1$.
\item No free charge and no free current in the bulk; the electric and magnetic problems are those of a
dielectric and a magnetisable medium.
\item Interface conditions inside the stack are ideal, with no slip, no debonding and no interface charge.
Imperfect interfaces are named as a limitation, not modelled.
\end{enumerate}

The second assumption is the one that decides whether the paper is publishable in a good journal or in a
mediocre one. If the two routes for the effective constants disagree by more than the thermal contribution to
attenuation, then the paper's central number is not measurable, and the honest response is to say so early
and to present the thermal effect as a bound. That decision is taken at week nine, on the evidence, not on
hope.

\section{The model}

\subsection{Constitutive relations}

With $\varepsilon_{ij}=\tfrac12(u_{i,j}+u_{j,i})$, $E_i=-\varphi_{,i}$ and $H_i=-\psi_{,i}$,
\begin{equation}
\bigl(1-l_{\rm nl}^2\nabla^2\bigr)t_{ij}=\sigma_{ij},\qquad
\sigma_{ij}=C_{ijkl}\bigl(1-l_g^2\nabla^2\bigr)\varepsilon_{kl}
-e_{kij}E_k-h_{kij}H_k-\beta_{ij}\theta ,
\end{equation}
\begin{equation}
D_i=e_{ikl}\varepsilon_{kl}+\epsilon_{ik}E_k+\alpha_{ik}H_k+p_i\theta ,
\end{equation}
\begin{equation}
B_i=h_{ikl}\varepsilon_{kl}+\alpha_{ik}E_k+\mu_{ik}H_k+m_i\theta ,
\end{equation}
with $h_{kij}$ piezomagnetic, $\alpha_{ik}$ magnetoelectric, $p_i$ pyroelectric and $m_i$ pyromagnetic. The
piezomagnetic constant is written with $h$ so that $q$ stays the carrier charge in the companion paper and
$k$ stays the wavenumber here. The two size lengths act on the elastic part of the mechanical stress and
nowhere else, so the electric and magnetic problems stay second order and need no extra condition at the
faces.

\subsection{Field equations}

\begin{equation}
\bigl(1-l_{\rm nl}^2\nabla^2\bigr)\rho\,\ddot u_i=\sigma_{ij,j},
\end{equation}
\begin{equation}
D_{i,i}=0,\qquad B_{i,i}=0,
\end{equation}
\begin{equation}
\Bigl(1+\tau_q\frac{\partial}{\partial t}\Bigr)
\Bigl[\rho c_p\dot\theta+T_0\beta_{ij}\dot\varepsilon_{ij}
+T_0p_i\dot E_i+T_0m_i\dot H_i\Bigr]=k_{\rm th}\nabla^2\theta .
\end{equation}
Five equations for five unknowns. The electrocaloric term $T_0p_i\dot E_i$ and the magnetocaloric term
$T_0m_i\dot H_i$ are the reciprocal partners of the pyroelectric and pyromagnetic terms in the constitutive
relations; dropping them while keeping $p_i$ and $m_i$ would break the reciprocity and is not done here. That
pairing is the only constitutive statement in this paper that is not already in the sources we cite, and it
is a textbook consequence rather than a discovery, so it goes in a sentence of the model section.

\subsection{Notation}

\begin{center}
\small
\begin{tabular}{p{2.4cm}p{6.7cm}p{3.4cm}}
\toprule
symbol & meaning & unit\\
\midrule
$u_x,u_z$ & displacements & m\\
$\varphi$, $\psi$ & electric and magnetic potentials & V, A\\
$\theta$ & temperature rise over $T_0$ & K\\
$k=k_r+ik_i$ & wavenumber; $k_i>0$ is decay along $+x$ & m$^{-1}$\\
$2h$ & plate thickness; faces at $z=\pm h$ & m\\
$C_{ijkl}$ & elastic moduli of the effective medium & Pa\\
$e_{kij}$ & piezoelectric constants & C\,m$^{-2}$\\
$h_{kij}$ & piezomagnetic constants & N\,A$^{-1}$m$^{-1}$\\
$\epsilon_{ik}$, $\mu_{ik}$ & dielectric permittivity, magnetic permeability & F\,m$^{-1}$, N\,A$^{-2}$\\
$\alpha_{ik}$ & effective magnetoelectric coefficient & s\,m$^{-1}$\\
$p_i$, $m_i$ & pyroelectric, pyromagnetic coefficients & C\,m$^{-2}$K$^{-1}$, N\,A$^{-1}$m$^{-1}$K$^{-1}$\\
$\beta_{ij}$ & thermal stress tensor $C_{ijkl}\alpha_{kl}$ & Pa\,K$^{-1}$\\
$l_{\rm nl}$, $l_g$ & nonlocal and gradient lengths, used as $l/h$ & m\\
$\rho c_p$, $k_{\rm th}$, $\kappa$ & heat capacity per volume, conductivity, diffusivity & J\,m$^{-3}$K$^{-1}$, W\,m$^{-1}$K$^{-1}$, m$^2$s$^{-1}$\\
$\tau_q$, $\tau_T$ & phase lag, $\tau_T=(2h)^2/(\pi^2\kappa)$ & s\\
$\omega$, $f$ & angular frequency, frequency & rad\,s$^{-1}$, Hz\\
\bottomrule
\end{tabular}
\end{center}

\subsection{Order of the system, and where this chapter stops}

The two operators raise the order of the problem in $z$, so the free faces need more than the classical
traction, electric and magnetic conditions, and the extra conditions must be consistent with the differential
form of the size effect if the answer is to mean anything. Which extra conditions to apply at $z=\pm h$,
and how to close the thermal condition there, is your problem to settle; this chapter stops before the
faces.

\section{Work plan}

Six steps. Each ends with a written comparison and, where a comparison is numerical, with the tolerance
stated in advance.

\paragraph{Step 1. Elastic plate, no fields} With $e=h=\alpha=0$, $\theta=0$ and both size lengths zero, the
model must return the classical Lamb spectrum of the laminated elastic plate and the Rayleigh velocity in
the short wavelength limit, to a relative $10^{-3}$ against the closed-form value computed from the same
$\rho$, $\lambda$ and $\mu$. Nothing else is computed until this passes.

\paragraph{Step 2. Electric and magnetic coupling without heat} Switch on $e_{kij}$, $h_{kij}$ and
$\alpha_{ik}$ and compare the branch structure with the closed-form multilayered results of Chen, Pan and
Chen \cite{Chen2007}. The comparison is of ordering and of the sign of the shift, since our constants come
from homogenisation and theirs from a different material set.

\paragraph{Step 3. Thermal coupling in the two limits} Switch on $\theta$ with $p_i$ and $m_i$ both retained,
and compare the isothermal and adiabatic limits. The coupled answer must lie between them to a width of order
the thermoelastic coupling. Then remove $m_i$ alone and report the change in attenuation; that number is how
much the paper's claim depends on the pyromagnetic term, and it should be stated whether it is large or
small.

\paragraph{Step 4. Relaxation switched on, and the two homogenisation routes} With $\tau_q$ finite, confirm the
saturation of attenuation at high frequency, the shift of the peak to $1/|\tau_T-\tau_q|$ and the rescaling
of the plateau, against a one-dimensional reduction of the same equations. Then repeat the calculation with
the effective constants from the second route, and draw the two sets of attenuation curves as a band. If the
band is wider than the thermal contribution, this is the paper's central negative result and it is reported
as such.

\paragraph{Step 5. Switching-off tests} Each coupling is switched off in two different ways, by setting its
constant to zero and by assembling a reduced model in which the term was never written. The two must agree to
$10^{-6}$ relative on $k_r$ and $k_i$ for every sampled mode: thermal off, electric off, magnetic off, size
effect off. The residual table goes in an appendix. A tolerance of $10^{-4}$ is a warning, not a pass.

\paragraph{Step 6. Size effect and the final sweep} Sweep $l_{\rm nl}/h$ and $l_g/h$ over the quoted range for
the composite, then produce the attenuation families for three thicknesses and both mode families. Report the
thermal contribution as a difference of attenuations so that a reader who ignores heat can still use the
mechanical curve.

\section{Figures, risks and timeline}

Four figures. First, dispersion of the elastic limit and of the full model on shared axes. Second,
attenuation of the two lowest branches against frequency times thickness, with three bands drawn: isothermal,
adiabatic, and the spread between the two homogenisation routes. Third, the band of thickness and phase lag
in which the thermal relaxation changes the peak by more than the linewidth, with the trust band of the
continuum heat model marked on the same axes. Fourth, the magnetoelectric contribution, obtained by removing
$\alpha_{ik}$, as a function of the layering ratio.

The risks. The first is resolution: the thermal contribution may be smaller than the spread between the two
homogenisation routes, in which case the paper reports a bound rather than a measurement of an effect. The
second is that the relaxation time of the composite is not one number, since the two phases have very
different conductivities, so an effective lag has to be defined and defended. The third is the size effect in
a laminated composite, where the relevant length is not that of either phase; carried as a sensitivity. The
fourth is the absence of a measured magnetoelectric coefficient for the exact stack we model, which we handle
by quoting the two routes and refusing to fit.

Timeline, sixteen weeks. Weeks one to three, material constants and homogenisation, since everything
downstream depends on them. Weeks four to five, Step 1. Weeks six to seven, Step 2. Weeks eight to nine,
Step 3 and the decision on whether the paper keeps its thermal claim, taken on the evidence with me.
Weeks ten to eleven, Step 4. Weeks twelve to thirteen, Step 5. Weeks fourteen to fifteen, Step 6 and figures.
Week sixteen, writing and the internal read-through.

\begin{thebibliography}{99}
\small
\bibitem{Chen2007} J.~Chen, E.~Pan, H.~Chen, \emph{Int.\ J.\ Solids Struct.} \textbf{44} (2007) 1073, \doi{10.1016/j.ijsolstr.2006.06.003}.
\bibitem{Xiao2016} D.~Xiao, Q.~Han, Y.~Liu, C.~Li, \emph{Compos.\ Struct.} \textbf{153} (2016) 704, \doi{10.1016/j.compstruct.2016.06.063}.
\bibitem{Kondaiah2013} P.~Kondaiah, K.~Shankar, N.~Ganesan, \emph{Coupled Syst.\ Mech.} \textbf{2} (2013) 1, \doi{10.12989/csm.2013.2.1.001}.
\bibitem{Tocci2021} G.~Tocci Monaco, N.~Fantuzzi, F.~Fabbrocino, R.~Luciano, \emph{Nanomaterials} \textbf{11} (2021) 87, \doi{10.3390/nano11010087}.
\bibitem{Kakar2015} R.~Kakar, S.~Kakar, \emph{Struct.\ Eng.\ Mech.} \textbf{53} (2015) 1201, \doi{10.12989/sem.2015.53.1201}.
\bibitem{Matlam2026} J.~L.~Matlam, G.~D.~Kedar, \emph{J.\ Therm.\ Stresses} \textbf{49} (2026) 1082, \doi{10.1080/01495739.2025.2612316}.
\bibitem{Lim2015} C.~W.~Lim, G.~Zhang, J.~N.~Reddy, \emph{J.\ Mech.\ Phys.\ Solids} \textbf{78} (2015) 298, \doi{10.1016/j.jmps.2015.02.001}.
\bibitem{Barretta2018} R.~Barretta, F.~Marotti de Sciarra, \emph{Int.\ J.\ Eng.\ Sci.} \textbf{130} (2018) 187, \doi{10.1016/j.ijengsci.2018.05.009}.
\end{thebibliography}
\end{document}
